\documentclass[12pt]{article}
\usepackage{amsmath,amssymb,amsthm,mathtools}
\usepackage{geometry}
\geometry{margin=1in}

\newtheorem{definition}{Definition}
\newtheorem{theorem}{Theorem}
\newtheorem{proposition}{Proposition}
\newtheorem{lemma}{Lemma}
\newtheorem{corollary}{Corollary}
\newtheorem{remark}{Remark}

\title{Mutual Liquidity Lock: A Formal Game-Theoretic Model}
\author{Working Draft}
\date{March 2026}

\begin{document}
\maketitle

%% ============================================================
%% 1. THE MLL GAME: FORMAL DEFINITION
%% ============================================================

\section{The MLL Game}

\begin{definition}[MLL Game]
The \textbf{Mutual Liquidity Lock (MLL) Game} is a discrete-time, infinite-horizon, two-player game $\Gamma = \langle \{A, B\}, \mathcal{A}, \mathcal{S}, u, \delta \rangle$ defined as follows.
\end{definition}

\subsection{Players and Time}

Two players $i \in \{A, B\}$. Periods $t = 0, 1, 2, \ldots$ with common discount factor $\delta \in (0,1)$.

\subsection{State Space}

The game state at time $t$ is a tuple:
\[
\omega_t = (s_t^A, s_t^B, \beta_t, \kappa_t, \gamma_t, \phi_t)
\]
where:
\begin{itemize}
    \item $s_t^i \geq 0$: player $i$'s \textit{share} (claim on the pool), with total pool $P_t = s_t^A + s_t^B$
    \item $\beta_t \in \{0, A, B\}$: bleeding target ($0$ = no bleeding active)
    \item $\kappa_t \in \{0, 1, \ldots, 6\}$: Russian Roulette invocation counter
    \item $\gamma_t \in \{0, 1, \ldots, \Theta\}$: cooldown timer ($\Theta = 10$ periods after RR invocation)
    \item $\phi_t \in \{\texttt{active}, \texttt{frozen}, \texttt{exited}\}$: game phase
\end{itemize}

\textbf{Initial state:} $\omega_0 = (s_0, s_0, 0, 0, 0, \texttt{active})$ where $s_0 > 0$ is each player's initial deposit.

\subsection{Action Space}

When $\phi_t = \texttt{active}$ and $\gamma_t = 0$ (no cooldown), each player $i$ simultaneously chooses:
\[
a_t^i \in \mathcal{A} = \{D, S, R, E\}
\]
\begin{itemize}
    \item $D$ (Deposit): contribute the agreed amount $d > 0$
    \item $S$ (Skip): do not contribute
    \item $R$ (Roulette): invoke the Russian Roulette mechanism
    \item $E$ (Exit): propose peaceful dissolution
\end{itemize}

During cooldown ($\gamma_t > 0$): $\mathcal{A} = \{D, S\}$ only.

When $\phi_t \in \{\texttt{frozen}, \texttt{exited}\}$: no actions available (terminal).

\subsection{Transition Rules}

\subsubsection{Yield Accrual}
Every period (regardless of actions), yield accrues at rate $y > 0$:
\[
\tilde{s}_t^i = s_t^i \cdot (1 + y)
\]
All subsequent computations in period $t$ use $\tilde{s}_t^i$.

\subsubsection{Deposit and Bleeding}

Given action profile $(a_t^A, a_t^B)$:

\textbf{Case 1: Both deposit} ($a_t^A = a_t^B = D$).
\[
s_{t+1}^i = \tilde{s}_t^i + d, \quad \beta_{t+1} = 0
\]

\textbf{Case 2: Player $i$ deposits, player $j$ skips} ($a_t^i = D, a_t^j = S$).

Bleeding activates against $j$:
\[
\Delta_t = r \cdot \tilde{s}_t^j
\]
\[
s_{t+1}^i = \tilde{s}_t^i + d + \Delta_t, \qquad s_{t+1}^j = \tilde{s}_t^j - \Delta_t = \tilde{s}_t^j(1 - r)
\]
\[
\beta_{t+1} = j
\]

\textbf{Case 3: Both skip} ($a_t^A = a_t^B = S$).

No bleeding transfer (symmetric default):
\[
s_{t+1}^i = \tilde{s}_t^i, \quad \beta_{t+1} = 0
\]

\begin{remark}
If both skip simultaneously, there is no asymmetry to trigger directional bleeding.
This creates a potential coordination problem analyzed in Section~\ref{sec:equilibrium}.
\end{remark}

\subsubsection{Peaceful Exit}

If both players choose $E$: $\phi_{t+1} = \texttt{exited}$. Each player receives their current share:
\[
u_{\text{exit}}^i = \tilde{s}_t^i
\]

If only one player chooses $E$, the proposal is recorded but no transition occurs. If both choose $E$ in the same period, exit executes.

\subsubsection{Russian Roulette Mechanism}
\label{sec:rr-def}

If any player chooses $R$ (and $\gamma_t = 0$):

\begin{enumerate}
    \item Counter increments: $\kappa_{t+1} = \kappa_t + 1$
    \item \textbf{Trigger check}: With probability $p(\kappa_{t+1}) = \frac{1}{7 - \kappa_{t+1}}$, funds freeze:
    \[
    \phi_{t+1} = \texttt{frozen}
    \]
    With probability $1 - p(\kappa_{t+1})$, the game continues with cooldown:
    \[
    \gamma_{t+1} = \Theta
    \]
    \item Cooldown decrements each period: $\gamma_{t+1} = \max(\gamma_t - 1, 0)$
\end{enumerate}

The probability sequence:
\[
p(k) = \frac{1}{7-k} \quad \text{for } k \in \{1, 2, 3, 4, 5, 6\}
\]

\begin{center}
\begin{tabular}{c|cccccc}
$k$ (invocation) & 1 & 2 & 3 & 4 & 5 & 6 \\
\hline
$p(k)$ & $\frac{1}{6}$ & $\frac{1}{5}$ & $\frac{1}{4}$ & $\frac{1}{3}$ & $\frac{1}{2}$ & $1$ \\
Cumulative & $\frac{1}{6}$ & $\frac{1}{3}$ & $\frac{1}{2}$ & $\frac{2}{3}$ & $\frac{5}{6}$ & $1$
\end{tabular}
\end{center}

\begin{remark}
This is equivalent to drawing without replacement from $\{1,2,3,4,5,6\}$.
The ``hidden $N$'' in the original design is isomorphic to this formulation,
but this version avoids the on-chain secrecy problem entirely: each invocation
requests fresh randomness (e.g., Chainlink VRF) with the appropriate probability.
\end{remark}

\subsection{Terminal Payoffs}

\subsubsection{Frozen Trust}
When $\phi = \texttt{frozen}$, principal is permanently locked. Each player receives a perpetual yield stream:
\[
u_{\text{frozen}}^i = \frac{y \cdot P_{\text{freeze}}}{2} \cdot \frac{1}{1 - \delta}
\]
where $P_{\text{freeze}} = s_{\text{freeze}}^A + s_{\text{freeze}}^B$ is the pool at freeze time.

Note: yield is split equally regardless of individual shares, as the ``frozen trust'' severs individual claims on principal. This is a design choice that maximizes the MAD deterrent.

\subsubsection{Peaceful Exit}
\[
u_{\text{exit}}^i = s_{\text{exit}}^i
\]
Each player receives their accumulated share (including any bleeding transfers).

\subsection{Per-Period Utility}

For player $i$ in an active period:
\[
u_t^i = y \cdot s_t^i - d \cdot \mathbf{1}[a_t^i = D]
\]

The total discounted utility from the game:
\[
U^i = \sum_{t=0}^{T-1} \delta^t \cdot u_t^i + \delta^T \cdot u_{\text{terminal}}^i
\]
where $T$ is the (possibly infinite) time of game termination.


%% ============================================================
%% 2. EQUILIBRIUM ANALYSIS
%% ============================================================

\section{Equilibrium Analysis}
\label{sec:equilibrium}

\subsection{The Cooperation Strategy Profile}

\begin{definition}[Cooperative Strategy $\sigma^*$]
For each player $i$:
\begin{enumerate}
    \item Play $D$ in every period where $\phi = \texttt{active}$ and the opponent has played $D$ in all prior periods (or is within a grace period $G$ after a skip).
    \item If opponent has played $S$ for more than $G$ consecutive periods, continue playing $D$ (collecting bleeding transfers).
    \item If bleeding has continued for $T_R > G$ periods with no resumption, play $R$ (initiate Russian Roulette).
    \item If opponent proposes $E$ and the pool split is acceptable (within threshold $\epsilon$ of proportional), play $E$.
\end{enumerate}
\end{definition}

\subsection{Deviation Analysis: Skip vs. Deposit}

\begin{theorem}[Cooperation Sustainability]
\label{thm:cooperation}
The cooperative strategy profile $\sigma^*$ constitutes a Subgame Perfect Equilibrium if and only if:
\[
\delta \geq \delta^* \equiv \frac{d}{d + r \cdot s_t^i}
\]
where $s_t^i$ is player $i$'s current share at the point of potential deviation.
\end{theorem}

\begin{proof}
Consider player $B$'s incentive to deviate from $D$ to $S$ at time $t$, given that $A$ follows $\sigma^*$.

\textbf{Payoff from cooperating} (continuing to play $D$):

Starting from share $s = s_t^B$ with yield rate $y$, the per-period net payoff is $y \cdot s_t^B - d$ (yield minus deposit). Under mutual cooperation the share grows:
\[
V_{\text{coop}}(s) = \sum_{\tau=0}^{\infty} \delta^\tau \left[ y(s + \tau d)(1+y)^\tau - d \right]
\]

For tractability, consider the marginal deviation incentive. In a single period:
\begin{itemize}
    \item Cooperate: net flow $= y \cdot s - d$, continuation share $= (s+d)(1+y)$
    \item Deviate (skip): net flow $= y \cdot s$, continuation share $= s(1+y)(1-r)$
\end{itemize}

The one-period gain from deviation: $+d$ (saved deposit).

The one-period cost of deviation (bleeding): in the next period, share is reduced by $r \cdot s(1+y)$. This compounds: the present value of the bleeding loss over $\tau$ periods before resuming is:
\[
\text{PV}_{\text{bleed}} = \sum_{\tau=1}^{\infty} \delta^\tau \left[ r \cdot s \cdot (1+y)^\tau \cdot (1-r)^{\tau-1} \right]
= \frac{\delta \cdot r \cdot s \cdot (1+y)}{1 - \delta(1+y)(1-r)}
\]

For deviation to be unprofitable:
\[
d \leq \frac{\delta \cdot r \cdot s \cdot (1+y)}{1 - \delta(1+y)(1-r)}
\]

When $y$ is small relative to $r$ and $\delta$, this simplifies to approximately:
\[
\delta \geq \frac{d}{d + r \cdot s}
\]

As $s$ grows over time (due to cumulative deposits), the RHS decreases, making cooperation \textit{easier} to sustain as the pool matures. $\qed$
\end{proof}

\begin{corollary}[Pool Maturity Effect]
The critical discount factor $\delta^*$ is decreasing in the pool size $P_t$. Thus, the MLL becomes more stable over time: the longer parties have cooperated, the harder it is to justify deviation.
\end{corollary}

\begin{remark}[Numerical Example]
With $d = 100$, $s = 10000$ (after several years of deposits), $r = 0.15$ (monthly bleeding, corresponding to $\approx 0.5\%$/day):
\[
\delta^* = \frac{100}{100 + 0.15 \times 10000} = \frac{100}{1600} = 0.0625
\]
Even extremely impatient agents (those who discount the future at 93.75\%) would find cooperation rational. The mechanism becomes very robust with accumulated stakes.
\end{remark}


\subsection{Russian Roulette Subgame Analysis}
\label{sec:rr-analysis}

When a player invokes Russian Roulette, we enter a finite extensive-form subgame. Let $V_{\text{RR}}^i(k)$ denote player $i$'s expected value when the counter is at $k$ and it is some player's turn to decide whether to invoke again (after cooldown).

\begin{definition}[Russian Roulette Subgame]
After invocation $k$ does not trigger ($k < 6$), both players enter a $\Theta$-period cooldown. At the end of cooldown, the game returns to the normal action space. At this point, either:
\begin{enumerate}
    \item Both resume cooperation (play $D$): the game returns to the normal repeated game.
    \item One or both skip: bleeding dynamics apply.
    \item Either player invokes $R$ again: counter goes to $k+1$.
    \item Both propose exit: peaceful dissolution.
\end{enumerate}
\end{definition}

\begin{proposition}[RR as Bargaining Catalyst]
\label{prop:rr-bargain}
In the Russian Roulette subgame with counter $\kappa = k$, if both players are risk-averse (concave utility over wealth), there exists a peaceful exit split $(\alpha, 1-\alpha)$ that both players prefer to continuing the RR process, provided:
\[
\frac{1}{7-k} > \frac{|s^A - s^B|}{P} \cdot c
\]
where $c$ is a constant depending on the curvature of the utility function.
\end{proposition}

\textit{Intuition}: As the trigger probability $p(k)$ increases with each invocation, the expected cost of continued brinkmanship rises. Risk-averse players will eventually prefer \textit{any} reasonable split to gambling on fund freeze. The key insight is that RR does not need to actually trigger to be effective --- the escalating probability drives convergence to a negotiated exit.

\begin{proposition}[RR Deterrence Threshold]
\label{prop:rr-deterrence}
A player will refrain from invoking Russian Roulette if and only if:
\[
V_{\text{status quo}}^i \geq (1 - p(k+1)) \cdot V_{\text{RR}}^i(k+1) + p(k+1) \cdot u_{\text{frozen}}^i
\]
where $V_{\text{status quo}}^i$ is the continuation value of the current (possibly bleeding) state.
\end{proposition}

\begin{remark}[When Does a Player Invoke RR?]
A player invokes RR when the status quo is sufficiently bad --- i.e., when bleeding has eroded their share enough that the gamble on fund freeze (which resets to equal yield sharing) becomes preferable. Specifically, if player $B$ is bleeding, $B$ invokes RR when:
\[
\frac{s_t^B}{P_t} < \frac{y}{2(1-\delta)} \cdot \frac{p(k+1)}{V_{\text{status quo}}^B}
\]
This creates a natural ``floor'' on how far bleeding can go before triggering escalation.
\end{remark}


%% ============================================================
%% 3. MECHANISM PROPERTIES
%% ============================================================

\section{Mechanism Properties}

\subsection{Individual Rationality}

\begin{theorem}[Ex-Ante Individual Rationality]
\label{thm:ir}
Participation in MLL is individually rational if:
\[
\frac{y \cdot s_0}{1 - \delta} - \frac{d}{1-\delta} + \frac{y \cdot d}{(1-\delta)^2} > s_0
\]
where the LHS is the expected discounted payoff from cooperation and the RHS is the outside option (keeping $s_0$).
\end{theorem}

\begin{proof}
Under mutual cooperation, player $i$'s share at time $t$ is:
\[
s_t^i = (s_0 + d \cdot t)(1+y)^t \approx s_0(1+y)^t + d \sum_{\tau=0}^{t-1}(1+y)^\tau
\]

The present value of yield minus deposits:
\[
V_{\text{coop}} = \sum_{t=0}^{\infty} \delta^t \left[ y \cdot s_t^i - d \right]
\]

For $V_{\text{coop}} > s_0$ (participation is better than autarky), we need the DeFi yield on the growing pool to compensate for the illiquidity cost. The condition simplifies to:
\[
y > (1 - \delta) \cdot \frac{d}{s_0 + \frac{d}{1-\delta}}
\]

With reasonable parameters ($\delta = 0.99$, $d = 100$, $s_0 = 1000$, $y = 0.005$ per period), this is satisfied. The yield on the locked pool must exceed the opportunity cost of illiquidity. $\qed$
\end{proof}

\subsection{Incentive Compatibility}

\begin{theorem}[Dominant Strategy in Mature Pool]
\label{thm:ic}
When $P_t$ is sufficiently large (specifically, when $r \cdot s_t^i > d / \delta$), playing $D$ is a dominant strategy regardless of the opponent's action.
\end{theorem}

\begin{proof}
If the opponent plays $D$: playing $D$ yields share growth of $d + y \cdot s_t^i$; playing $S$ yields share growth of $y \cdot s_t^i$ but triggers bleeding risk if opponent retaliates.

If the opponent plays $S$: playing $D$ yields $d + y \cdot s_t^i + r \cdot s_t^j$ (deposit plus bleeding bonus); playing $S$ yields only $y \cdot s_t^i$ (no bleeding, no growth).

In both cases, the gain from $D$ over $S$ is at least $d + r \cdot s_t^j - d = r \cdot s_t^j > 0$ when accounting for the bleeding transfer.

More precisely, the net gain from playing $D$ vs. $S$ when the opponent plays $S$:
\[
\underbrace{r \cdot s_t^j}_{\text{bleeding bonus}} - \underbrace{d}_{\text{deposit cost}} + \underbrace{\delta \cdot r \cdot (s_t^i + d)(1+y)}_{\text{PV of avoiding being bled next period}}
\]

When $r \cdot s_t^j > d$ (pool share exceeds deposit), $D$ dominates $S$ even myopically. With discounting, the condition relaxes further. $\qed$
\end{proof}

\subsection{Renegotiation-Proofness}
\label{sec:reneg}

\begin{definition}[Renegotiation-Proofness]
A strategy profile is \textit{renegotiation-proof} if there is no period $t$ and no alternative continuation strategy profile that both players strictly prefer to the prescribed continuation.
\end{definition}

\begin{proposition}[Conditional Renegotiation-Proofness]
\label{prop:reneg}
The cooperative equilibrium $\sigma^*$ is renegotiation-proof if and only if in every punishment phase (bleeding), the punishing player weakly prefers to continue punishment over any mutually beneficial renegotiation.
\end{proposition}

\textit{Analysis}: This is the most subtle property. Consider the situation where $B$ has defaulted and is bleeding. Both players know that:
\begin{enumerate}
    \item $A$ is gaining: $A$ receives bleeding transfers and would prefer this to continue.
    \item $B$ is losing: $B$ would prefer to renegotiate.
    \item If they renegotiate to peaceful exit, $A$ still gets a favorable split (due to accumulated bleeding transfers).
\end{enumerate}

The key question is whether $A$ prefers continued bleeding or immediate exit. If $A$'s continuation value from bleeding exceeds the exit value, $A$ will refuse renegotiation, making the punishment self-enforcing.

\[
A \text{ refuses renegotiation} \iff \sum_{\tau=1}^{\infty} \delta^\tau \cdot r \cdot s_t^B (1-r)^{\tau-1}(1+y)^\tau > 0
\]

This is \textit{always} satisfied when $s_t^B > 0$ and $r < 1$. Therefore:

\begin{corollary}
The bleeding punishment is inherently renegotiation-proof from the punisher's perspective: the punishing player always prefers to continue collecting bleeding transfers over stopping. The threat of Russian Roulette by the punished party serves as the only constraint on indefinite bleeding.
\end{corollary}

This creates a natural dynamic equilibrium: bleeding continues until the punished party's threat of Russian Roulette becomes credible enough to induce negotiation. The ``negotiation zone'' emerges endogenously.


%% ============================================================
%% 4. PARAMETER CALIBRATION
%% ============================================================

\section{Parameter Calibration}

\subsection{Bleeding Rate $r$}

The bleeding rate must satisfy three constraints simultaneously:

\begin{enumerate}
    \item \textbf{Deterrence}: $r$ must be high enough that the expected cost of skipping exceeds the deposit amount:
    \[
    r > \frac{d(1-\delta)}{s_t^i \cdot \delta}
    \]

    \item \textbf{Recovery allowance}: $r$ must be low enough that a temporarily incapacitated player can recover meaningful share:
    \[
    (1-r)^{G} > \alpha_{\min}
    \]
    where $G$ is the grace period (in periods) and $\alpha_{\min}$ is the minimum acceptable share retention (e.g., 80\%).

    \item \textbf{RR trigger calibration}: $r$ should be calibrated so that the punished party's RR invocation threshold is reached after a meaningful but not excessive number of periods:
    \[
    T_{\text{RR}} = \frac{\ln(\theta_{\text{RR}})}{\ln(1-r)} \quad \text{periods until RR becomes rational}
    \]
    where $\theta_{\text{RR}}$ is the share ratio at which RR invocation becomes the best response.
\end{enumerate}

\begin{proposition}[Optimal Bleeding Rate]
For the baseline parameterization ($d = 100$, $s_0 = 1000$, $\delta = 0.99$, $y = 0.005$), the optimal bleeding rate lies in:
\[
r^* \in [0.003, 0.008] \text{ per day} \quad \Leftrightarrow \quad [0.09, 0.22] \text{ per month}
\]

The original design's $r = 0.005$/day ($\approx 0.14$/month) falls within this range.
\end{proposition}


\subsection{Russian Roulette Probability Sequence}

\begin{proposition}[Optimality of the $1/(7-k)$ Sequence]
The probability sequence $p(k) = 1/(7-k)$ for $k \in \{1,\ldots,6\}$ is equivalent to sampling without replacement from a uniform distribution on $\{1,\ldots,6\}$. This has two desirable properties:
\begin{enumerate}
    \item \textbf{Uniform marginal}: each invocation is equally likely to be the ``fatal'' one ex ante ($\Pr[\text{trigger at exactly invocation } k] = 1/6$ for all $k$).
    \item \textbf{Monotone hazard rate}: $p(k)$ is strictly increasing in $k$, ensuring escalation pressure mounts with each invocation.
\end{enumerate}
\end{proposition}

\begin{remark}[Design Alternative: Geometric Distribution]
An alternative would be i.i.d. draws with fixed probability $p$ each round (geometric distribution). This has a constant hazard rate, which means escalation pressure does not increase. The without-replacement design is strictly preferable for brinkmanship because it guarantees termination in finite time and creates accelerating urgency.
\end{remark}


%% ============================================================
%% 5. EXTENSIONS AND OPEN QUESTIONS
%% ============================================================

\section{Extensions and Open Questions}

\subsection{Asymmetric Deposits}

When players have different incomes, equal deposits create asymmetric exit costs. Let $w^A, w^B$ be incomes with $w^A > w^B$. With equal deposit $d$:
\[
\text{Relative cost: } \frac{d}{w^A} < \frac{d}{w^B}
\]

Player $B$ faces higher relative exit cost. This can be addressed by:
\begin{enumerate}
    \item \textbf{Proportional deposits}: $d^i = \alpha \cdot w^i$ (requires income verification --- difficult on-chain)
    \item \textbf{TCP discovery}: the exploratory protocol described in the original design
    \item \textbf{Endogenous asymmetry}: allow different deposit amounts; adjust bleeding rates proportionally to maintain symmetric incentives
\end{enumerate}

\subsection{Multi-Party Extension (M-of-N)}

Extending to $N > 2$ players with $M$-of-$N$ signing threshold introduces:
\begin{itemize}
    \item Coalition formation: subsets of players may collude
    \item Voting on exit: majority vs. unanimity
    \item Directed bleeding: who bleeds when one player defaults?
\end{itemize}

This connects to the broader mechanism design literature on team production (Holmstr\"om, 1982) and requires a separate treatment.

\subsection{Dead Man's Switch}

To address the key liveness concern (one party losing access to their key), define:
\begin{definition}[Dead Man's Switch]
If player $i$ fails to interact with the contract for $T_{\text{dead}}$ consecutive periods, player $j$ gains unilateral withdrawal rights over the entire pool.
\end{definition}

The parameter $T_{\text{dead}}$ must be set large enough to avoid false triggers but small enough to be practically useful. Suggested range: $T_{\text{dead}} \in [60, 180]$ days.

\subsection{Smart Contract Completeness vs. Real-World Incompleteness}

MLL is a \textit{complete} contract in the Hart-Moore sense: all state transitions are deterministically specified by the code. This avoids the renegotiation problem of incomplete contracts (Hart \& Moore, 1988) but creates rigidity. The ``escape valve'' is the peaceful exit mechanism, which requires bilateral agreement but allows any arbitrary wealth transfer between parties.


\end{document}
